\documentclass[12pt]{article}
\usepackage[utf8]{inputenc}
\usepackage[T1]{fontenc}
\usepackage{hyperref}
\title{Dai Zhuang Language Resources for K-12 Curriculum Development}
\author{Compiled by Claude Sonnet 5.5 (claude-sonnet-5-5) for Andr\'as Kornai}
\date{}
\begin{document}
\maketitle

\section{Introduction}
Dai Zhuang (ISO 639-3 zhd; Glottocode daiz1235; also called Wenma, Thu Lao and Tuliao) is a Tai language of the Zhuang macrolanguage. It is spoken in Yunnan, mainly in Wenshan Prefecture, and in northern Vietnam. Searches found general reference pages only. They found no dictionaries, grammars or school materials specific to this variety. Wikipedia lists zhd and tyl (Thu Lao) as duplicate codes, so check which code a resource uses.

\section{Primary Educational Resources}
No primers, textbooks or school programmes specific to Dai Zhuang were found. Standard Zhuang has its own literature and school use, but it is based on a different variety, and its suitability for Dai Zhuang was not assessed.

\section{Linguistic and Reference Materials}
\begin{itemize}
\item Wikipedia article on Dai Zhuang, with the dialect division (Johnson 2011) and geography: \url{https://en.wikipedia.org/wiki/Dai_Zhuang_language}
\item Wikipedia overview of the Zhuang languages: \url{https://en.wikipedia.org/wiki/Zhuang_languages}
\item Wikipedia article on Standard Zhuang, including the Sawndip script: \url{https://en.wikipedia.org/wiki/Standard_Zhuang}
\item Omniglot page on the Zhuang writing systems: \url{https://omniglot.com/writing/zhuang.htm}
\end{itemize}

\section{Religious and Community Resources}
No scripture translations or community organisations specific to Dai Zhuang were found. Global Recordings Network lists the variety, but the search result says it has no recordings.

\section{Digital and Audio Resources}
\begin{itemize}
\item Global Recordings Network page for Zhuang, Dai: Central Wenshan, which lists no recordings: \url{https://www.globalrecordings.net/en/language/25592}
\item Global Recordings Network page for the zhd code: \url{https://glorecnet-website.k.globalrecordings.net/en/language/zhd}
\end{itemize}

\section{Curriculum Development Context}
Dai Zhuang is used as a first language by all adults but not by all young people. A curriculum would need documentation specific to this variety. Standard Zhuang materials and the Sawndip tradition could serve as a starting point, but whether they suit Dai Zhuang needs community input. The search results contained no information on this.

\end{document}